\capitulocapa{images/cap-3.png}{Como compartilhar estado?}
\label{chap:compartilhar-estado}

\section{A solução apareceu no review}

\begin{historia}
Depois da tentativa do capítulo anterior, o pull request tinha propriedades
criadas às pressas, um segundo estado no Perfil e um comentário nada animador:

``Funciona, mas o Header virou dono de uma preferência usada pelo sistema
inteiro.''

Durante o review, alguém fez a pergunta certa:

``E se o tema subir para o Layout, que já envolve o Header, a Sidebar e as
páginas?''

A mudança não exigia uma biblioteca. Exigia colocar a fonte de verdade num
lugar compatível com quem precisava dela.
\end{historia}

No Capítulo~2, percebemos que o tema deixou de ser local. Agora vamos executar o
primeiro movimento importante de arquitetura em React: elevar o estado para o
ancestral comum.

Essa solução é simples, explícita e frequentemente suficiente. Também possui um
limite. Vamos atravessá-lo sem fingir que ele não existe.

\section{O problema: componentes irmãos precisam concordar}

O DevPanel possui uma estrutura parecida com esta:

\begin{lstlisting}
<AppLayout>
  <Header />
  <Sidebar />
  <MainContent>
    <ProfilePage />
  </MainContent>
</AppLayout>
\end{lstlisting}

O Header altera o tema. A Sidebar precisa exibir o tema atual. O Perfil também
precisará conhecê-lo. Nenhum desses componentes é ancestral dos outros.

Manter o estado no Header significa tentar enviar informação de um irmão para
outro. React não oferece um caminho lateral para isso, e criar um caminho
escondido normalmente resulta em DOM, variáveis de módulo ou eventos globais
difíceis de rastrear.

O caminho previsível é subir até o primeiro componente que contém todos os
consumidores e, a partir dele, distribuir o valor para baixo.

\imagemcapitulo
  {ch03-fluxo-fonte-unica}
  {O AppLayout se torna a fonte única e distribui o tema aos consumidores.}
  {fluxo-fonte-unica}

\section{Como normalmente resolvemos}

Uma reação comum é colocar o estado no componente mais alto possível, muitas
vezes no \codigo{App}. Isso funciona, mas pode transformar a raiz numa coleção
de estados sem relação entre si: tema, modal, formulário, filtros, notificações
e qualquer outra informação que alguém precisou compartilhar.

O objetivo não é subir tudo para o topo. É subir apenas até o ancestral comum
mais próximo.

No DevPanel, esse lugar é \codigo{AppLayout}. Ele representa a região que contém
o Header, a Sidebar e as páginas autenticadas. Uma tela de login fora desse
Layout não precisa receber a preferência neste momento.

\begin{decisao}
Elevar estado não significa tornar estado global. Significa mover a fonte de
verdade até a menor fronteira capaz de atender todos os consumidores atuais.
\end{decisao}

\section{Uma fonte de verdade no Layout}

Movemos o tipo, o valor e a ação para \codigo{AppLayout}:

\begin{lstlisting}[caption={O Layout passa a ser o dono do tema.},label={lst:theme-layout}]
import { useState } from 'react'

export type Theme = 'light' | 'dark'

export function AppLayout() {
  const [theme, setTheme] = useState<Theme>('light')

  function toggleTheme() {
    setTheme((current) => (current === 'light' ? 'dark' : 'light'))
  }

  return (
    <div className={`app app--${theme}`}>
      <Header theme={theme} onToggleTheme={toggleTheme} />
      <Sidebar theme={theme} />
      <MainContent />
    </div>
  )
}
\end{lstlisting}

Agora o estado não depende da presença do Header. Enquanto o Layout existir, o
tema existe. Header e Sidebar recebem a mesma fotografia do valor.

\section{Implementação nos consumidores}

O Header deixou de criar estado. Seu contrato mostra exatamente o que ele usa:

\begin{lstlisting}[caption={O Header recebe valor e ação.}]
type HeaderProps = {
  theme: Theme
  onToggleTheme: () => void
}

export function Header({ theme, onToggleTheme }: HeaderProps) {
  return (
    <header>
      <span>Tema: {theme === 'light' ? 'claro' : 'escuro'}</span>
      <ThemeToggle theme={theme} onToggle={onToggleTheme} />
    </header>
  )
}
\end{lstlisting}

A Sidebar precisa apenas ler:

\begin{lstlisting}[caption={A Sidebar não recebe uma ação que não utiliza.}]
type SidebarProps = {
  theme: Theme
}

export function Sidebar({ theme }: SidebarProps) {
  return (
    <aside aria-label="Navegação principal">
      <span>Tema atual: {theme === 'light' ? 'claro' : 'escuro'}</span>
      <nav>{/* links do DevPanel */}</nav>
    </aside>
  )
}
\end{lstlisting}

Essa diferença é importante. Receber o valor não deveria conceder permissão para
alterá-lo. Passamos a ação somente a quem oferece aquela interação.

\section{Entendendo o fluxo}

Quando a pessoa seleciona ``Usar tema escuro'', a sequência é esta:

\begin{enumerate}
  \item \codigo{ThemeToggle} chama \codigo{onToggle};
  \item o Header chama a ação recebida do Layout;
  \item \codigo{toggleTheme} atualiza o estado no Layout;
  \item o Layout renderiza novamente com \codigo{theme} igual a
    \codigo{'dark'};
  \item Header e Sidebar recebem o novo valor;
  \item os dois representam o mesmo tema.
\end{enumerate}

Não existe sincronização entre Header e Sidebar. Existe uma fonte de verdade e
dois consumidores. Essa distinção elimina a pergunta ``como manter os dois
estados iguais?'': só há um estado.

\subsection{Dados descem}

O valor segue do proprietário para os filhos por propriedades. Cada componente
consegue ler seu contrato e descobrir de onde a informação chega.

\subsection{Eventos sobem}

O botão não modifica diretamente o Layout. Ele comunica uma intenção por uma
função recebida. A função é executada no componente que possui o estado.

Esse fluxo unidirecional é previsível: dados descem, eventos sobem, uma nova
renderização distribui o resultado.

\section{Limitações: quando as props só estão de passagem}

A solução começa a incomodar quando o Perfil, renderizado dentro de
\codigo{MainContent}, também precisa do tema:

\begin{lstlisting}
<AppLayout>
  <MainContent theme={theme}>
    <RouteContent theme={theme}>
      <ProfilePage theme={theme} />
    </RouteContent>
  </MainContent>
</AppLayout>
\end{lstlisting}

Os componentes intermediários talvez não usem o tema. Ainda assim, seus
contratos precisam aceitá-lo apenas para continuar o caminho. Isso é
\emph{props drilling}: propriedades atravessando camadas intermediárias para
alcançar um consumidor distante.

\imagemcapitulo
  {ch03-arvore-props-drilling}
  {Componentes intermediários recebem e repassam uma propriedade que não usam.}
  {arvore-props-drilling}

Props drilling não é automaticamente um erro. Duas ou três camadas numa árvore
estável podem ser mais claras que adicionar outra abstração. O problema aparece
quando o custo de transportar a informação começa a dominar os componentes.

Alguns sinais:

\begin{itemize}
  \item componentes recebem propriedades que não leem;
  \item mover uma tela exige reconstruir todo o caminho;
  \item várias propriedades percorrem sempre as mesmas camadas;
  \item componentes intermediários mudam por causa de uma demanda distante;
  \item o contrato deixa de representar a responsabilidade real do componente.
\end{itemize}

\section{Uma solução melhor: Context API}

Antes de buscar outro mecanismo, podemos tentar composição. Em vez de
\codigo{MainContent} criar internamente toda a árvore, o Layout pode receber o
conteúdo já montado. Outra opção é aproximar o componente que precisa do estado.

Essas mudanças funcionam quando a estrutura é o verdadeiro problema. Quando a
informação é transversal por natureza --- como tema, idioma ou sessão ---,
reorganizar a árvore apenas para transportar props pode piorar o desenho.

É aqui que Context API começa a fazer sentido. Não porque props sejam ruins, mas
porque existe um valor transversal, muitos consumidores e intermediários sem
interesse nele.

Context permite que um Provider disponibilize um valor a qualquer descendente,
sem obrigar todas as camadas intermediárias a recebê-lo. A fonte de verdade
continua única; muda apenas o mecanismo de distribuição.

\begin{lstlisting}[caption={O contexto descreve o contrato compartilhado.}]
import { createContext, useContext } from 'react'

type ThemeContextValue = {
  theme: Theme
  toggleTheme: () => void
}

const ThemeContext = createContext<ThemeContextValue | null>(null)

export function useTheme() {
  const value = useContext(ThemeContext)

  if (!value) {
    throw new Error('useTheme precisa estar dentro de ThemeProvider')
  }

  return value
}
\end{lstlisting}

O valor inicial é \codigo{null} porque não existe tema válido fora do Provider.
O hook \codigo{useTheme} centraliza a leitura e produz um erro claro quando a
árvore foi configurada incorretamente.

\section{O Provider mantém a fonte de verdade}

\begin{lstlisting}[caption={O Provider oferece valor e ação aos descendentes.}]
type ThemeProviderProps = {
  children: ReactNode
}

export function ThemeProvider({ children }: ThemeProviderProps) {
  const [theme, setTheme] = useState<Theme>('light')

  function toggleTheme() {
    setTheme((current) => (current === 'light' ? 'dark' : 'light'))
  }

  return (
    <ThemeContext value={{ theme, toggleTheme }}>
      {children}
    </ThemeContext>
  )
}
\end{lstlisting}

No React 19, o próprio contexto pode ser usado como Provider. O DevPanel envolve
sua área autenticada uma única vez:

\begin{lstlisting}
<ThemeProvider>
  <AppLayout />
</ThemeProvider>
\end{lstlisting}

Agora Header, Sidebar e Perfil leem diretamente aquilo que utilizam:

\begin{lstlisting}
export function Sidebar() {
  const { theme } = useTheme()

  return <aside>Tema atual: {theme}</aside>
}
\end{lstlisting}

\codigo{AppLayout} e \codigo{MainContent} deixam de transportar props. Eles não
precisam conhecer uma preferência que não usam.

\section{Exercício: remova o transporte}

\begin{tarefa}
Migre o tema compartilhado para \codigo{ThemeProvider}. Faça Header, Sidebar e
Perfil usarem \codigo{useTheme} e remova de AppLayout e MainContent todas as
propriedades que existiam apenas para transportar o tema.
\end{tarefa}

Depois de implementar, registre:

\begin{enumerate}
  \item quais propriedades puderam ser removidas;
  \item quais componentes agora dependem de \codigo{ThemeProvider};
  \item onde a fonte de verdade continua localizada;
  \item o que aconteceria se um consumidor ficasse fora do Provider;
  \item por que Context resolveu distribuição, mas não definiu as ações.
\end{enumerate}

\espacoresposta{9}

Compare a versão final com a implementação por props. A nova versão deve reduzir
o transporte sem esconder quem realmente consome o tema.

\begin{resumo}
\begin{itemize}
  \item Componentes irmãos compartilham estado por meio de um ancestral comum.
  \item O estado deve subir apenas até a menor fronteira que atende seus
    consumidores.
  \item Uma fonte de verdade elimina a necessidade de sincronizar cópias.
  \item Valores descem por propriedades e eventos sobem por callbacks.
  \item Props drilling importa quando intermediários participam sem usar os
    dados.
  \item Context é justificável quando o valor é transversal e a árvore vira
    apenas um canal de transporte.
  \item O Provider mantém a fonte de verdade; o contexto distribui seu contrato.
  \item Um hook próprio torna o consumo consistente e falha de forma clara.
\end{itemize}
\end{resumo}

Compartilhar estado não começa com Zustand. Começa escolhendo quem é responsável
por ele e tornando o fluxo visível para o time. No próximo capítulo, veremos o
que acontece quando Context deixa de ser um canal pequeno e vira um depósito.
